\section{Before the Christian Rite: Ancient Contexts}

Christian exorcism did not arise in a world without discourse about hostile
spiritual agency. The witnesses below supply context, not a genealogy.
Similar acts can rest on incompatible accounts of God, spirits, illness,
ritual efficacy, and authorized ministry; resemblance does not prove literary
dependence.

\subsection{Israel's Scriptures: affliction without a ritual manual}

The Hebrew Bible resists being flattened into a handbook. The harmful spirit
in 1 Samuel 16 belongs to a narrative of kingship and judgment; David's music
gives Saul relief, but the narrator supplies neither adjuration nor ministerial
office. Tobit 6--8 joins Raphael's direction, Tobias and Sarah's prayer,
marriage, material means, healing, and angelic action under providence. Its
Greek recensions and Latin reception do not always arrange or phrase the
episode identically. The history must therefore identify its textual form.
Neither narrative is a primitive Roman rite or a general authorization to
imitate its material action.

\subsection{Second Temple plurality}

In R.~H. Charles's English translation, \emph{Jubilees} 10.1--11 depicts
destructive spirits afflicting Noah's descendants and Noah petitioning God.
After Mastema asks that some remain under him, God permits one tenth to remain
and consigns nine tenths to condemnation. This is a literary episode in
rewritten Scripture, not evidence of one standard Jewish ceremony; Charles's
translation is not an original-language control.

Josephus presents himself as having seen Eleazar perform before Vespasian and
a military audience. The account says that God gave Solomon skill against
demons; Eleazar uses a ring containing a root, mentions Solomon while reciting
incantations attributed to him, and commands the departing demon to overturn a
water-filled cup or basin. The locus is \emph{Antiquities} 8.45--49 in modern
Niese numbering and VIII.2.5 in Whiston's division. It establishes Josephus's
late-first-century literary testimony and Solomonic framing, not the mechanism,
actual Solomonic authorship, or an unbroken route into the Christian rite.

Together these witnesses show why ``the Jewish background'' is too singular:
genre, means, agent, theology, and purpose differ.

\subsection{Comparison without a borrowed-rite thesis}

The draft does not convert a broad list of Mesopotamian, Egyptian, Greek, or
Roman practices into a supposed background rite. Such comparison requires
genre-, date-, and edition-specific witnesses. Without a shared formula,
traceable route, or explicit ancient claim of borrowing, resemblance remains
context rather than descent.
